\scanpage{223}
\mspage{- 59 -}
n\textsuperscript{o} 215

\uline{structure\unsure{s} associée à} $S$, $(f,f')$ \uline{un homomorphisme de} $E$ \uline{dans un fût} $E'$, $E^T$ \uline{et} ${E'}^T$ \uline{les fûts d'espèce} $T$ \uline{associés à} $E$ \uline{et} $E'$, $(g,f')$ \uline{l'homomorphisme associé à} $(f,f')$ \uline{du tronc sous-jacent de} $E^T$ \uline{dans celui de} ${E'}^T$ ; $(g,f')$ \uline{est alors un homomorphisme de} $E^T$ \uline{dans} ${E'}^T$.

Il suffira manifestement de montrer ce qui suit : soient $D$ et $D'$ des débitages de $E$ et $E'$ qui soient associés par $(f,f')$ (le rédacteur s'excuse de l'emploi multiple du mot \char34 associé\char34{} ; il a eu la flemme de tout recommencer avec un autre mot), $D^T$ et ${D'}^T$ les débitages de $E^T$ \unsure{\uline{et}|et} ${E'}^T$ associés à $D$, $D'$ respectivement ; $D^T$ et ${D'}^T$ sont alors associés par $(g,f')$. Soit $U$ le domaine de $D$, et soit $x$ un point de $U$ ; posons $x'=f'(x)$. Si $y\in U$, $D$ définit un isomorphisme $J_x$ \unsure{de} la fibre $E_x$ de $x$ dans $E$ sur celle de $y$, et $D'$ un isomorphisme de la fibre $E'_{x'}$ de $x'$ dans $E'$ sur celle de $y'=f'(y)$. Il résulte du fait que $D$, $D'$ sont associés par $(f,f')$ que l'on a
\[
f\circ J_y=J'_{y'}\circ f.
\]
En effet, si $\xi\in T_x$, soit $s_\xi$ la section appartenant à $D$ telle que $s_\xi(x)=\xi$, d'où $J_y(\xi)=s_\xi(y)$. Il y a par hypothèse une section $s'$ appartenant à $D'$ telle que $f\circ s_\xi$ soit la restriction de $s'\circ f'$ à $U$. On a $(f\circ J_y)(\xi)=s'(y')=J'_{y'}(s'(x'))=J'_{y'}((f\circ J_x)(\xi))=(J'_{y'}\circ f)(\xi)$, puisque $J_x$ est l'automorphisme identique. Soient $f_x$, $f_y$ les restrictions de $f$ à $E_x$, $E_y$, $g_x$ et $g_y$ les restrictions de $g$ aux fibres $E_x^T$, $E_y^T$ de $x$, $y$ dans $E^T$ ; $g_x$ et $g_y$ sont donc les restrictions à $E_x^T$, $E_y^T$ des extensions canoniques de $f_x$, $f_y$ aux échelons (des échelles construites sur $E_x$, $E_y$) dont font partie $E_x^T$, $E_y^T$. Soient $J_y^T$ l'isomorphisme de $E_x^T$ sur $E_y^T$ défini par le débitage $D^T$ et ${J'}^{T}_{y'}$ celui de ${E'}^{T}_{x'}$ sur ${E'}^{T}_{y'}$ défini par le débitage ${D'}^T$ ; $J_y^T$ (resp. : ${J'}^{T}_{y'}$) est donc la restriction à $E_x^T$
% archive sheet number: 222
